\chapter{Conclusions and Future Work}
\label{chap:conclusoes-e-trabalhos-futuros}

This work proposed the HoTHP, a variant of the Transformer Hawkes Process that replaces the trigonometric kernel of the RoTHP with a hyperbolic kernel that has exponential decay. The motivation was to improve length extrapolation in temporal point processes: to train on short sequences and keep performance on longer sequences.

\section{Contributions}

The main contribution is the introduction of the exponential recency bias into the temporal attention mechanism. Unlike the oscillatory kernel of the RoTHP, whose attention score can grow for temporal distances not seen in training, the hyperbolic kernel of the HoTHP guarantees a monotonic decay. For any temporal distance, including values larger than those observed during training, the attention score is bounded and tends to zero. This allows the model to generalize to longer sequences without degradation.

Besides the architectural proposal, the experimental analysis revealed two observations that we consider relevant for the field:

\begin{enumerate}
    \item The exponential bias acts exclusively on the attention mechanism, which builds the hidden state from the past events. The intensity function is a separate layer with free learnable parameters. This means that the HoTHP can model both self-exciting and self-inhibiting processes. What the bias imposes is not the shape of the intensity, but the shape of the relevance: recent events receive more weight in the construction of the internal representation.

    \item The visualization of the learned intensities showed that the Amazon dataset exhibits a self-inhibiting pattern (the intensity grows between events and falls when an event occurs), while Stackoverflow exhibits an almost constant intensity (no temporal memory). The HoTHP wins on the first and ties on the second, confirming that the relevant factor is the presence of short-range temporal memory, regardless of the direction (excitation or inhibition).
\end{enumerate}

\section{Main results}

On the synthetic data generated by Hawkes processes with an exponential kernel, the HoTHP obtains the best NLL at all extrapolation factors ($2\times$, $5\times$, and $10\times$) in both decay regimes, surpassing all the baselines (NHP, THP, and RoTHP); the advantage is statistically significant (paired, one-sided Wilcoxon test) in all cases, with the slow decay at $10\times$ against the NHP being the narrowest-margin comparison. This result was expected, since the bias of the model is perfectly aligned with the generating process.

On the real data, the results vary by dataset. The HoTHP obtains the best performance on Amazon ($\Delta\text{NLL} = -0{,}074$ in the $5\times$ extrapolation, against $-0{,}020$ of the RoTHP), ties on Stackoverflow and Taxi, and fails on Retweet with raw data due to the extreme temporal scale of the dataset. The ablation with temporal normalization revealed that the Retweet problem was primarily one of numerical scale: with normalized timestamps, the HoTHP obtains an NLL $5\times$ of $-0{,}651 \pm 0{,}013$, the best result among all the models and conditions tested, with almost zero variance between seeds. The analysis of these results points to two conditions necessary for the bias to be beneficial: the process needs to have short-range temporal memory, and the scale of the intervals between events needs to be in a range that allows the calibration of the parameter $\theta'$.

From the computational point of view, the HoTHP is between $1{,}05\times$ and $2{,}24\times$ slower in training, depending on the length of the sequences. This additional cost is paid only once and is justified by the benefit in extrapolation: the trained model produces stable results on sequences up to $5\times$ longer, while the RoTHP can collapse numerically in this same scenario.

\section{Advantages and disadvantages of Transformers for point processes}

The experiments also give a more balanced picture of what the Transformer brings to point process modeling. On the positive side, attention removes the sequential bottleneck of the recurrent models: all the events are processed in parallel, which makes training faster on long sequences. Our own timing measurements show this clearly, since the recurrent NHP is the slowest model of all, between $1{,}3\times$ and $1{,}5\times$ slower than the HoTHP. Attention also relates any two events directly, without carrying information through a single hidden state, and it offers a natural place to inject a temporal inductive bias, which is exactly what the HoTHP does through its kernel. On the negative side, attention has no native notion of continuous time: by itself it is invariant to the order of the events, so the temporal structure has to be added through the positional encoding, and the choice of this encoding is decisive, as the gap between the THP, the RoTHP, and the HoTHP shows. The attention cost is also quadratic in the length of the sequence, against the linear cost of a recurrent model. And, on real data, the recurrent NHP still obtains the best likelihood on several datasets, because its continuous-time formulation models the density between events natively, something that attention only approximates.

Taken together, the results suggest that there is no single best model, and that the choice depends on the goal. When the priority is to train on a large volume of data or on long sequences, the parallelism of the Transformer is the decisive advantage. When the goal is to predict the type of the next event, the HoTHP obtains the best accuracy, surpassing even the recurrent baseline on Amazon and Stackoverflow. When the priority is stable performance on sequences much longer than those seen in training, the recency bias of the HoTHP is what keeps the model from degrading. And when the goal is only the likelihood on a real dataset with an unknown structure, the recurrent NHP tends to win, unless the process has the kind of decaying memory that the HoTHP is built for, as in the synthetic experiments and under extrapolation, where the HoTHP obtains the best NLL as well. The contribution of this work should be read in this light: not as a model that dominates every metric, but as one that makes the Transformer respect the temporal decay structure of the Hawkes process and, with it, extrapolate in a stable way.

\section{Limitations}

The work has some limitations that should be considered:
\begin{enumerate}
    \item The evaluation on real data is limited to four datasets. A broader evaluation, including datasets from domains such as seismology, finance, and computer networks, would strengthen the conclusions about when the bias is beneficial.

    \item The HoTHP does not have an internal mechanism for temporal normalization. The scale of the timestamps is handled only by the shift to zero ($\tilde{t}_i = t_i - t_1$), which preserves the differences but does not control the absolute magnitude. Datasets with very large intervals (such as Retweet) cause numerical saturation in the kernel.

    \item The parameter $\theta'$ controls a single global decay rate. Processes with multiple temporal scales may require multiple decay rates.
\end{enumerate}

\section{Future work}
